% GENERATED FILE. Edit canonical lessons, then run npm run generate:books.
% canonical-source-hash: fnv1a64:f3af0989c1a91b68
% canonical-lessons: TA-C433-cikaram, TA-C433-malaiyeru, TA-C433-erimalai, TA-C433-panipparai, TA-C433-camaveli

\chapter{Peaks and Plains}
\label{ch:ta-peaks-and-plains}
\hlchaptermodality{433}

\begin{chapteropening}
\textbf{By the end of this chapter:} \emph{I can talk about peaks, volcanoes, glaciers and plains, and trekking.}

Complete the last lesson of chapter 433: I can talk about peaks, volcanoes, glaciers and plains, and trekking.
\end{chapteropening}

\section[\texorpdfstring{\ta{சிகரம்} (cikaram)}{சிகரம் (cikaram)}]{\ta{சிகரம்} (cikaram) — a peak; a summit}

\label{lesson:TA-C433-cikaram}



\begin{quote}
Before the new one: say the Tamil for a water well, then the Tamil for to slosh.
\end{quote}

\subsection*{You'll want to know: \ta{சிகரம்}}
\textbf{\ta{சிகரம்}} — \emph{cikaram} — \textquotedblleft{}a peak; a summit\textquotedblright{}.

\begin{quote}
\ta{சிகரத்தில்} \ta{எப்போதும்} \ta{பனி} \ta{இருக்கிறது}.
\end{quote}

\begin{quote}
\emph{cikarattil eppōtum paṉi irukkiṟatu.}
\end{quote}

\begin{quote}
There is always snow on the peak.
\end{quote}

A -\ta{அம்} noun: \ta{சிகரத்தில்}, on the peak.

\begin{cousinweb}
\textbf{\ta{சிகரம்}} (\emph{cikaram}) is borrowed from Sanskrit \emph{śikhara}, \textquotedblleft{}a point, a peak\textquotedblright{} (CDIAL 12435).
\end{cousinweb}

\begin{grammarlens}[title={What you've built}]
One more word for getting around.
\end{grammarlens}

\subsection*{Your turn}
\begin{itemize}
\raggedright
  \item \emph{Say it:} \emph{cikaram}
  \item \emph{Say it:} \emph{cikaram}, once more
  \item \emph{Say it:} \emph{cikarattil eppōtum paṉi irukkiṟatu.} — There is always snow on the peak.
  \item \emph{Recall:} say the Tamil for a water well, then the Tamil for to slosh, then say \emph{cikaram} again
\end{itemize}

\begin{tcolorbox}[breakable,colback=teal!4,colframe=teal!35!black,title={Before you move on}]
What does \ta{சிகரம்} mean? (\textquotedblleft{}A peak; a summit\textquotedblright{}.) Say it once more.
\end{tcolorbox}
\section[\texorpdfstring{\ta{மலையேறு} (malaiyēṟu)}{மலையேறு (malaiyēṟu)}]{\ta{மலையேறு} (malaiyēṟu) — to trek; to climb mountains}

\label{lesson:TA-C433-malaiyeru}



\begin{quote}
Before the new one: say the Tamil for to slosh, then the Tamil for a peak.
\end{quote}

\subsection*{You'll want to know: \ta{மலையேறு}}
\textbf{\ta{மலையேறு}} — \emph{malaiyēṟu} — \textquotedblleft{}to trek; to climb mountains\textquotedblright{}.

\begin{quote}
\ta{விடுமுறையில்} \ta{நான்} \ta{நண்பர்களுடன்} \ta{மலையேறினேன்}.
\end{quote}

\begin{quote}
\emph{viṭumuṟaiyil nāṉ naṇparkaḷuṭaṉ malaiyēṟiṉēṉ.}
\end{quote}

\begin{quote}
In the holidays I went trekking with friends.
\end{quote}

Past (I …): \ta{மலையேறினேன்}, \emph{malaiyēṟiṉēṉ}.

\ta{மலை}, a hill, + \ta{ஏறு}, to climb, written as one word.

\begin{cousinweb}
\textbf{\ta{மலையேறு}} (\emph{malaiyēṟu}) is built from \ta{மலை}, \textquotedblleft{}a hill, a mountain\textquotedblright{} (DEDR 4742), and \ta{ஏறு}, \textquotedblleft{}to climb\textquotedblright{} (DEDR 916).
\end{cousinweb}

\begin{grammarlens}[title={What you've built}]
One more word for getting around.
\end{grammarlens}

\subsection*{Your turn}
\begin{itemize}
\raggedright
  \item \emph{Say it:} \emph{malaiyēṟu}
  \item \emph{Say it:} \emph{malaiyēṟu}, once more
  \item \emph{Say it:} \emph{viṭumuṟaiyil nāṉ naṇparkaḷuṭaṉ malaiyēṟiṉēṉ.} — In the holidays I went trekking with friends.
  \item \emph{Recall:} say the Tamil for to slosh, then the Tamil for a peak, then say \emph{malaiyēṟu} again
\end{itemize}

\begin{tcolorbox}[breakable,colback=teal!4,colframe=teal!35!black,title={Before you move on}]
What does \ta{மலையேறு} mean? (\textquotedblleft{}To trek; to climb mountains\textquotedblright{}.) Say it once more.
\end{tcolorbox}
\section[\texorpdfstring{\ta{எரிமலை} (erimalai)}{எரிமலை (erimalai)}]{\ta{எரிமலை} (erimalai) — a volcano}

\label{lesson:TA-C433-erimalai}



\begin{quote}
Before the new one: say the Tamil for a peak, then the Tamil for to trek.
\end{quote}

\subsection*{You'll want to know: \ta{எரிமலை}}
\textbf{\ta{எரிமலை}} — \emph{erimalai} — \textquotedblleft{}a volcano\textquotedblright{}.

\begin{quote}
\ta{அந்தத்} \ta{தீவில்} \ta{ஒரு} \ta{எரிமலை} \ta{இருக்கிறது}.
\end{quote}

\begin{quote}
\emph{antat tīvil oru erimalai irukkiṟatu.}
\end{quote}

\begin{quote}
There is a volcano on that island.
\end{quote}

\ta{எரி}, to burn, + \ta{மலை}, a mountain: a burning mountain.

\begin{cousinweb}
\textbf{\ta{எரிமலை}} (\emph{erimalai}) is built from \ta{எரி}, \textquotedblleft{}to burn\textquotedblright{} (DEDR 811), and \ta{மலை}, \textquotedblleft{}a mountain\textquotedblright{} (DEDR 4742).
\end{cousinweb}

\begin{grammarlens}[title={What you've built}]
One more word for getting around.
\end{grammarlens}

\subsection*{Your turn}
\begin{itemize}
\raggedright
  \item \emph{Say it:} \emph{erimalai}
  \item \emph{Say it:} \emph{erimalai}, once more
  \item \emph{Say it:} \emph{antat tīvil oru erimalai irukkiṟatu.} — There is a volcano on that island.
  \item \emph{Recall:} say the Tamil for a peak, then the Tamil for to trek, then say \emph{erimalai} again
\end{itemize}

\begin{tcolorbox}[breakable,colback=teal!4,colframe=teal!35!black,title={Before you move on}]
What does \ta{எரிமலை} mean? (\textquotedblleft{}A volcano\textquotedblright{}.) Say it once more.
\end{tcolorbox}
\section[\texorpdfstring{\ta{பனிப்பாறை} (paṉippāṟai)}{பனிப்பாறை (paṉippāṟai)}]{\ta{பனிப்பாறை} (paṉippāṟai) — a glacier; an iceberg}

\label{lesson:TA-C433-panipparai}



\begin{quote}
Before the new one: say the Tamil for to trek, then the Tamil for a volcano.
\end{quote}

\subsection*{You'll want to know: \ta{பனிப்பாறை}}
\textbf{\ta{பனிப்பாறை}} — \emph{paṉippāṟai} — \textquotedblleft{}a glacier; an iceberg\textquotedblright{}.

\begin{quote}
\ta{பனிப்பாறைகள்} \ta{மெதுவாக} \ta{உருகுகின்றன}.
\end{quote}

\begin{quote}
\emph{paṉippāṟaikaḷ metuvāka urukukiṉṟaṉa.}
\end{quote}

\begin{quote}
The glaciers are slowly melting.
\end{quote}

\ta{பனி}, snow, + \ta{பாறை}, a rock, with the \ta{ப்} doubled.

\begin{cousinweb}
\textbf{\ta{பனிப்பாறை}} (\emph{paṉippāṟai}) is built from \ta{பனி}, \textquotedblleft{}dew, snow\textquotedblright{} (DEDR 4035), and \ta{பாறை}, \textquotedblleft{}a rock\textquotedblright{} (DEDR 4121).
\end{cousinweb}

\begin{grammarlens}[title={What you've built}]
One more word for getting around.
\end{grammarlens}

\subsection*{Your turn}
\begin{itemize}
\raggedright
  \item \emph{Say it:} \emph{paṉippāṟai}
  \item \emph{Say it:} \emph{paṉippāṟai}, once more
  \item \emph{Say it:} \emph{paṉippāṟaikaḷ metuvāka urukukiṉṟaṉa.} — The glaciers are slowly melting.
  \item \emph{Recall:} say the Tamil for to trek, then the Tamil for a volcano, then say \emph{paṉippāṟai} again
\end{itemize}

\begin{tcolorbox}[breakable,colback=teal!4,colframe=teal!35!black,title={Before you move on}]
What does \ta{பனிப்பாறை} mean? (\textquotedblleft{}A glacier; an iceberg\textquotedblright{}.) Say it once more.
\end{tcolorbox}
\section[\texorpdfstring{\ta{சமவெளி} (camaveḷi)}{சமவெளி (camaveḷi)}]{\ta{சமவெளி} (camaveḷi) — a plain; flat land}

\label{lesson:TA-C433-camaveli}



\begin{quote}
Before the new one: say the Tamil for a volcano, then the Tamil for a glacier.
\end{quote}

\subsection*{You'll want to know: \ta{சமவெளி}}
\textbf{\ta{சமவெளி}} — \emph{camaveḷi} — \textquotedblleft{}a plain; flat land\textquotedblright{}.

\begin{quote}
\ta{மலையிலிருந்து} \ta{கீழே} \ta{பரந்த} \ta{சமவெளி} \ta{தெரிந்தது}.
\end{quote}

\begin{quote}
\emph{malaiyiliruntu kīḻē paranta camaveḷi terintatu.}
\end{quote}

\begin{quote}
From the hill a wide plain could be seen below.
\end{quote}

\ta{சம}, level, + \ta{வெளி}, an open space.

\begin{cousinweb}
\textbf{\ta{சமவெளி}} (\emph{camaveḷi}) is built from \emph{cama}, \textquotedblleft{}level\textquotedblright{}, borrowed from Sanskrit \emph{sama}, \textquotedblleft{}equal, level\textquotedblright{} (CDIAL 13173), and \ta{வெளி}, \textquotedblleft{}an open space\textquotedblright{} (DEDR 5498).
\end{cousinweb}

\begin{grammarlens}[title={What you've built}]
One more word for getting around.
\end{grammarlens}

\subsection*{Your turn}
\begin{itemize}
\raggedright
  \item \emph{Say it:} \emph{camaveḷi}
  \item \emph{Say it:} \emph{camaveḷi}, once more
  \item \emph{Say it:} \emph{malaiyiliruntu kīḻē paranta camaveḷi terintatu.} — From the hill a wide plain could be seen below.
  \item \emph{Recall:} say the Tamil for a volcano, then the Tamil for a glacier, then say \emph{camaveḷi} again
\end{itemize}

\begin{tcolorbox}[breakable,colback=teal!4,colframe=teal!35!black,title={Before you move on}]
What does \ta{சமவெளி} mean? (\textquotedblleft{}A plain; flat land\textquotedblright{}.) Say it once more.
\end{tcolorbox}
